\chapter{Conclusion}


As a result of all these experiments and discussions, the following conclusions can be drawn:



\section{RQ1}

The images generated by the StyleGAN-C generator at 256x256 pixels can be measured in three ways. First, FID, which measures how close the distributions between two samples are, achieves a value of 46.08, significantly improving the other SOTA results. Second, CAS, which calculates how well the images synthesized by conditional generators are classified, results in 0.923, meaning the classifier correctly classifies 0.923 of the generated images. Finally, HITL metrics indicate that experts can only diagnose 0.50 of the images correctly, with most diagnoses made for pituitary tumors. \textbf{Therefore, it can be concluded that the generated images are highly useful for the machine and the classifier, but for the expert, only some images are useful.} This can be attributed to the low quality of the dataset, limiting robust learning for all classes.

\section{RQ2}

In the project, StyleGAN-C was used for inversion via optimization and an encoder, where the metrics were 0.838 for accuracy (CNN), 0.765 for LPIPS, and 0.832 for SVM (e4e), achieving non-trivial results. The inversion model is more effective but computationally more expensive. Additionally, StyleGAN-C has the ability to generate counterfactual videos as a structural validator. On the other hand, when using StyleGAN-C for class balancing and CNN as the classifier, an accuracy of 0.960 is reached, far surpassing the inversion method. \textbf{Therefore, it can be stated that StyleGAN-C performs better as an image generation model complementing classification models than for classification via inversion, although inversion can also serve as a classifier.}

\section{RQ3}

When testing with Grad-CAM with experts, it was found that there is no difference in the real images, as the expert diagnoses all cases correctly. However, for the fake images, confidence when making the diagnosis increased from 1.44 to 3.0, and accuracy improved from 0.44 to 0.56, thus \textbf{Grad-CAM is particularly useful for synthetic images.}

\section{Hypothesis}


\begin{itemize}

    \item \textbf{H1:} \textbf{is partially supported}, as the average diagnostic utility for the pituitary class is 3.8, matching that of real images. However, for meningioma and glioma classes, the results are 1.67 and 1.33, respectively. Only half of the generated images were diagnostically useful.

    
    \item \textbf{H2A:} \textbf{is fully supported}, where both experiments (inversion with optimization and with encoder) achieved accuracy metrics of 0.838 and 0.832, respectively.
    
    \item \textbf{H2B:} \textbf{is fully supported}, where the metrics statistically improved from a single-run baseline of 0.934 accuracy with only real images to 0.96 using an ensemble and class balancing with StyleGAN-C images. The model that includes images generated by 
    StyleGAN is statistically more reliable than the model without them in ensembles, increasing from 0.953 to 0.96 with a p-value of 0.0065.
    
    \item \textbf{H3:} \textbf{is partially supported}, where confidence increased from 1.44 to 3.0, and accuracy from 0.44 to 0.56 in synthetic images with Grad-CAM, compared to when Grad-CAM is not used. For real images, a comparison cannot be made since the accuracy is 100.
\end{itemize}

\section*{General Conclusion \footnote{Contact: junweihemai1@gmail.com}}

This thesis addresses three key aspects. It proposes an end-to-end framework that encompasses the training of an image generator and the explanation of its classification. Based on this, it can be concluded that:

\begin{enumerate}
    \item An \textbf{image generator is trained to synthesize conditional brain tumor images}, accurately selecting tumor types with superior quality compared to the SOTA.
    
    
    \item  \textbf{A CNN classifier with strict leakage control} is proposed, achieving metrics of \textbf{0.96}, surpassing SOTA.
    
    \item For the first time, a \textbf{pilot HITL integration} is proposed to review synthetic images and heatmap results generated by Grad-CAM in brain tumor classification, as well as the use of the \textbf{CAS metric} to measure the utility of synthetic conditional images in the classifier.
\end{enumerate}

\section{Future Work}

First, a key direction is to pursue \textbf{stronger external validation}. All experiments in this thesis were conducted on a single public dataset, split into complete and reduced variants under a strict patient-level separation into internal and external partitions. As a result, the robustness of the proposed framework under domain shift remains an open question. Future work should therefore evaluate the methodology on multi-centre, multi-scanner cohorts and explicitly study performance under distribution shifts. Such external validation is essential to assess whether the gains in calibration, explainability, and accuracy observed here transfer to more heterogeneous clinical environments. Such external validation is essential to establish the external validity of the proposed framework and to assess whether the observed gains in calibration, explainability, and accuracy persist under real world deployment conditions.
\\

Second, the current framework operates at the \textbf{2D slice level} and does not explicitly exploit volumetric or multi-planar information. An important extension would be to design tri-planar (axial/sagittal/coronal) fusion models or fully 3D architectures that integrate information across planes and slices for each tumor. Moving from a purely 2D slice-wise setup to tri-planar or 3D modelling would provide the generative model with richer contextual information, potentially enabling a more faithful representation of tumor extent and shape for glioma and meningioma.\\

Third, the results of Experiment~2C and the HITL ablations suggest that \textbf{careful data curation and meningioma under representation} still limit performance. Future work should thus prioritize the collection of additional, high-quality meningioma cases and more balanced test sets. The HITL results, however, indicate that a sizable fraction of synthetic images are either non diagnostic or potentially misleading, especially for glioma and meningioma. A natural extension of the proposed framework is therefore to introduce an explicit curation stage, in which synthetic samples are filtered and, when feasible, expert review. Such a filtering mechanism would aim to retain only high-quality, diagnostically plausible images for augmentation, thereby reducing the risk of propagating generator artefacts into the discriminative model.\\


Fourth, Experiment~1 shows that truncation can significantly alter both tumor extent and anatomy, motivating research on \textbf{more disentangled and controllable latent spaces}. Ideally, tumor morphology should be editable while preserving the surrounding brain structure. Possible avenues include explicit factorization of tumor versus background factors, spatially localized style modulation, or the use of segmentation masks during training to constrain where generative changes are allowed to occur. These ideas connect naturally with the volumetric and tri-planar extensions outlined above.\\


Fifth, while the present work focuses on unimodal MRI-based classification, a natural direction for future research is to extend the framework towards genuine \textbf{multimodal fusion} by incorporating structured clinical metadata such as age, sex, symptom duration, or tumor grade. These variables could be integrated either as additional inputs to the discriminative classifier or as conditioning information for the generative model. Such a multimodal design would align more closely with real-world diagnostic workflows, in which radiological findings are interpreted in conjunction with patient history, and could further improve both predictive performance and the clinical relevance of the generated counterfactuals. In parallel, richer and more standardized patient-level metadata would enable the training of large language models (LLMs) or related modules that reason over textual and clinical information in combination with MRI-derived features. Furthermore, the rapid advancement of large multimodal models such as GPT-4o and Gemini opens the possibility of leveraging these foundation models for medical image analysis. While their current capabilities in controlled class-conditional generation and patient-level evaluation remain limited compared to task-specific architectures like StyleGAN-C, their potential as complementary tools for clinical decision support warrants investigation in future work.\\

Finally, as a long-term goal, the framework could be integrated into hospitals, helping as a \textbf{clinical decision support system}. In this setup, the models would provide calibrated probabilities, counterfactual explanations and curated synthetic examples, along with explainable AI (XAI) to aid in understanding model predictions, with the final decision remaining with the clinician. Evaluating the impact of such a system on reading time, diagnostic confidence, and inter-observer variability would be a key step towards clinical adoption.\\

The present study opens several promising research directions. First, it suggests that generative classification under strict patient-level leakage control is a viable and principled alternative for studying robustness and generalization in medical imaging, particularly in data-scarce and clinically realistic settings. Second, it motivates the development of functional evaluation metrics for synthetic data, such as CAS, that move beyond perceptual realism and explicitly assess alignment with downstream diagnostic tasks. Finally, the proposed use of generative counterfactuals points toward a new class of explainable AI methods grounded in anatomically plausible alternative hypotheses, enabling structured reasoning and deeper model interrogation. Together, these directions position generative models not only as auxiliary tools for performance improvement, but as foundational components for robust, interpretable, and methodologically sound clinical decision-support systems.


\section{Applications to the Real World}

\begin{itemize}
    \item \textbf{Image generation for students}: Due to privacy concerns and the lack of access to public data, these generators can be used to synthesize images for educational purposes, allowing students to learn and practice diagnostic skills without compromising patient confidentiality. 
    
    \item \textbf{Use of XAI for difficult images}: It can be seen that experts can diagnose with high accuracy, but when blurry or low-quality images (in this case, the fakes) are used, accuracy decreases. XAI can also aid in the segmentation of difficult regions, providing additional insights to improve diagnosis in such cases.

    \item \textbf{Early disease detection}: The system can assist in the early detection of various medical conditions, improving patient outcomes by enabling timely interventions in hospital settings.

    \item \textbf{Access to diagnostics in remote areas}: Where specialists may be scarce, an AI-powered decision support system can assist healthcare providers through telemedicine. This system aids clinicians with diagnostic tools and expert insights, without replacing the medic. By leveraging AI, it improves access to medical care in regions with limited resources.

    \item \textbf{Extension to classification of other diseases:} 
    The model can be extended to the classification of other medical conditions by leveraging similar computer vision techniques. For example, this approach could be adapted to classify other types of tumors (e.g., lung, breast, or prostate cancer) by training on respective medical image datasets.

\end{itemize}


